\documentclass{beamer}

\usetheme{Madrid}

\title[Introduction to LaTeX]{Introduction to LaTeX}
\author[Saarvik Singh]{Saarvik Singh Suryavanshi}
\institute[MNIT Jaipur]{
	MNIT Jaipur \\
	Technical Writing \\
	Course Coordinator: Neeta Nain
}
\date{\today}

\setbeamertemplate{navigation symbols}{}

\setbeamertemplate{footline}{
	\leavevmode
	\hbox{
		\begin{beamercolorbox}[
			wd=0.34\paperwidth,
			ht=2.5ex,
			dp=1ex,
			left
			]{author in head/foot}
			\hspace{1em}
			Saarvik Singh Suryavanshi
		\end{beamercolorbox}
		
		\begin{beamercolorbox}[
			wd=0.46\paperwidth,
			ht=2.5ex,
			dp=1ex,
			center
			]{title in head/foot}
			Introduction to LaTeX
		\end{beamercolorbox}
		
		\begin{beamercolorbox}[
			wd=0.20\paperwidth,
			ht=2.5ex,
			dp=1ex,
			right
			]{date in head/foot}
			\insertframenumber/\inserttotalframenumber
			\hspace{1em}
		\end{beamercolorbox}
	}
}

\begin{document}
	
	\begin{frame}
		\titlepage
	\end{frame}
	
	\section{Introduction}
	
	\begin{frame}{What is LaTeX?}
		
		\begin{itemize}
			\item LaTeX is a document preparation system
			based on TeX.
			
			\vspace{0.3cm}
			
			\item It is commonly used for academic papers,
			reports, and technical documents.
			
			\vspace{0.3cm}
			
			\item It provides consistent formatting and
			good support for mathematical expressions.
			
			\vspace{0.3cm}
			
			\item It uses commands to define the structure
			and appearance of a document.
		\end{itemize}
		
	\end{frame}
	
	\section{Basic Features}
	
	\begin{frame}{Basic Features of LaTeX}
		
		\begin{block}{Document Structure}
			A LaTeX document contains a preamble and
			a document environment.
		\end{block}
		
		\begin{block}{Mathematical Expressions}
			LaTeX makes it easy to write equations.
			
			\[
			E = mc^2
			\]
		\end{block}
		
		\begin{block}{Other Features}
			\begin{itemize}
				\item Tables and figures
				\item Automatic numbering
				\item Sections and chapters
				\item References and citations
			\end{itemize}
		\end{block}
		
	\end{frame}
	
	\section{Conclusion}
	
	\begin{frame}{Conclusion}
		
		\begin{itemize}
			\item LaTeX is useful for preparing technical
			and academic documents.
			
			\vspace{0.3cm}
			
			\item It helps maintain consistent formatting.
			
			\vspace{0.3cm}
			
			\item Its mathematical and document-structuring
			features make it useful for students and researchers.
		\end{itemize}
		
		\vspace{0.8cm}
		
		\begin{center}
			\Large\textbf{Thank You!}
		\end{center}
		
	\end{frame}
	
\end{document}